\section{Results}
\label{sec:results}
% Every table cites its results file with \src{...} and marks each value F, P or not run.

\subsection{An inpainting-evaluator case study}
Saved synthetic probes of BraTS \texttt{inpainting} 0.0.12 identify a two-spatial-axis SSIM convention and PSNR clipping that breaks exact rank equivalence with MSE.
They motivate an explicit implementation contract, not a claim of independent anatomical sensitivity.
The supplement retains both dependency-version records and all numerical probes; no generator or organizer-server validation is implied.
% Sources: results/brats_panel_checks_tm1.1.2.json and brats_panel_checks_tm1.9.0.json.

\subsection{Brain MRI (QTIM)}
\label{sec:r-brain}
All brain-panel findings in this section are provisional.
Table~\ref{tab:brain} reports the saved random-region, both-hemisphere profile, with source two-way intervals.
The MSE position of the dataset mean is below the mirror position, while the corresponding ordering changes under TM-like distance.
The difference is an observed metric profile under the recorded alignment, not evidence that a population mean preserves more of a person's anatomy.

The tests also distinguish separate requirements.
MSE establishes that the repeat is better than the dataset mean, incorrect detail, and noise in the recorded directional tests.
It does not establish that the repeat is better than its attenuated version.
The attenuated-repeat MSE position is $-0.014$ with interval $[-0.028,-0.004]$\statuspre, where the original repeat is the zero anchor.
For the repeat-tolerance lDDT-like measure, the corresponding position is $-0.084$ with interval $[-0.108,-0.065]$\statuspre.
These findings show a possible attenuation preference while leaving the registration gate unresolved.
They do not make failure of T2 a general theorem about those metrics.
% Numerical source: donor_ladder_brain_metrics.json, views/random_patches/LR/ladder_R.

\begin{table*}[t]
\centering
\caption{\textbf{Provisional brain audit profile.} $R_{\rm rep}$: repeat zero and mean of the recorded unrelated donors one; lower is closer.
Cells contain means and source 95\% intervals, with family and position resampling.
Non-repeat rows involve 221 targets from 168 families; the repeat anchor involves 23 targets.
Residual correlation at the exact dataset mean is undefined.
Source: \src{results/donor_ladder_brain_metrics.json}.}
\label{tab:brain}
\small
\begin{tabular}{@{}lcccc@{}}
\toprule
Metric & Dataset mean & Mirror & Incorrect detail & Noise \\
\midrule
MSE & \shortstack{0.413\statuspre\\{\scriptsize [0.367, 0.443]}} & \shortstack{0.717\statuspre\\{\scriptsize [0.622, 0.792]}} & \shortstack{1.006\statuspre\\{\scriptsize [0.992, 1.020]}} & \shortstack{0.921\statuspre\\{\scriptsize [0.893, 0.944]}} \\
SSIM & \shortstack{0.738\statuspre\\{\scriptsize [0.680, 0.784]}} & \shortstack{0.823\statuspre\\{\scriptsize [0.755, 0.873]}} & \shortstack{1.023\statuspre\\{\scriptsize [1.014, 1.036]}} & \shortstack{1.132\statuspre\\{\scriptsize [1.088, 1.194]}} \\
Residual correlation & \textit{undefined} & \shortstack{0.752\statuspre\\{\scriptsize [0.688, 0.808]}} & \shortstack{1.003\statuspre\\{\scriptsize [0.996, 1.009]}} & \shortstack{1.003\statuspre\\{\scriptsize [1.000, 1.006]}} \\
TM-like & \shortstack{0.723\statuspre\\{\scriptsize [0.642, 0.778]}} & \shortstack{0.667\statuspre\\{\scriptsize [0.530, 0.761]}} & \shortstack{0.997\statuspre\\{\scriptsize [0.987, 1.007]}} & \shortstack{1.075\statuspre\\{\scriptsize [1.041, 1.121]}} \\
lDDT-like, repeat tolerance & \shortstack{0.355\statuspre\\{\scriptsize [0.254, 0.430]}} & \shortstack{0.890\statuspre\\{\scriptsize [0.843, 0.925]}} & \shortstack{1.050\statuspre\\{\scriptsize [1.030, 1.080]}} & \shortstack{2.073\statuspre\\{\scriptsize [1.924, 2.273]}} \\
\bottomrule

\end{tabular}
\end{table*}

HACA3+ stage 1 provides a separate contrast sensitivity, not a completed identity endpoint.
Across its 23 paired scans, mean absolute deep-region session shift changes from $0.121$ to $0.016$ standardized intensity units\statuspre.
The parietal counterpart changes from $0.059$ to $0.033$\statuspre.
The source records a pending visual verdict, and neither measurement establishes preservation of subject-specific detail.
% Numerical source: donor_ladder_brain_haca3_stage1.json, session_shift_z.

\subsection{Fundus photographs (PAPILA, GRAPE)}
\label{sec:r-fundus}
The frozen bilateral-branch aggregate supplies the P8/P9 frame-only profile in Table~\ref{tab:fundus}.
GRAPE has 239 target-eye units from 134 people in the stranger comparison; the repeat comparison has 171 units from 97 people\statuspre.
PAPILA has 488 target eyes from 244 people and no other-visit rung\statuspre.
The dataset mean is closer than the fellow eye under squared error in both saved summaries.
The context-selected mean of look-alikes can outperform the unconditional mean in this representation.
These findings are provisional and do not establish the identity of the inferred detail.

\begin{table*}[t]
\centering
\caption{\textbf{Provisional frame-only fundus profile.}
Mean-relative squared-error ratio $Q$, with source person-bootstrap intervals; lower is closer.
These P8/P9 map-compatible summaries differ from the older P7 support and from the window-level metric panel.
PAPILA has no repeat observation.
Source: \src{results/donor_ladder_2d_facefmt.json} at the frozen bilateral snapshot, copied without change into the paper inputs.}
\label{tab:fundus}
\small
\begin{tabular}{@{}lcccc@{}}
\toprule
Dataset & Repeat & Fellow eye & Context mean & Stranger \\
\midrule
GRAPE P8 & \shortstack{0.489\statuspre\\{\scriptsize [0.404, 0.585]}} & \shortstack{1.330\statuspre\\{\scriptsize [1.224, 1.437]}} & \shortstack{0.856\statuspre\\{\scriptsize [0.832, 0.880]}} & \shortstack{2.012\statuspre\\{\scriptsize [1.897, 2.140]}} \\
PAPILA P9 & \textit{unavailable} & \shortstack{1.467\statuspre\\{\scriptsize [1.422, 1.512]}} & \shortstack{0.935\statuspre\\{\scriptsize [0.914, 0.955]}} & \shortstack{1.992\statuspre\\{\scriptsize [1.941, 2.049]}} \\
\bottomrule

\end{tabular}
\end{table*}

The Gaussian maps in Figure~\ref{fig:atlas} retain these preprocessing labels.
High split-half correlation cannot establish anatomical validity if a shared preprocessing artifact is reproducible too.
The completed v1 metric registry and the running newer extension remain separate.
Superseded support-sensitive block searches supply no primary result.

\subsection{Face photographs (CelebA, IMDb-WIKI)}
\label{sec:r-face}
Table~\ref{tab:face} reports the available v6f tier-1 profiles on the self-anchored scale.
The repeat photograph has substantial pixel error despite preserving person membership.
The dataset mean is closer than the repeat under MSE in both datasets, while a local context-selected donor is closer still.
That pixel preference does not identify the donor as the target person.

For CelebA, the selected donor's local MSE position is $0.399$\statusfinal, while its whole-image ArcFace donor-reference position is $0.983$\statusfinal.
The corresponding IMDb-WIKI values are $0.397$ and $0.984$\statusfinal.
These are different measurements at different spatial scopes.
Their disagreement shows that local photometric matching can select a donor whose global identity representation remains near the unrelated reference.
It does not measure identity retention by a completed composite reconstruction, and no ArcFace value is assigned to an unembedded mean.

\begin{table*}[t]
\centering
\caption{\textbf{Face metric and donor-identity profiles.} Recorded $R_{\rm self}$ point estimates; target self zero, mean of unrelated donors one.
The source retains family/position bootstrap intervals; full records accompany the preparation manifest.
The approved frame has 921 CelebA targets and 678 IMDb-WIKI targets, respectively clustered into 921 and 453 units.
ArcFace compares whole target and donor images; the pixel metrics compare local windows.
Source: \src{results/face_ladder_metrics.json}.}
\label{tab:face}
\small
\begin{tabular}{@{}llcccc@{}}
\toprule
Dataset & Metric & Retest & Bilateral & Dataset mean & Context donor \\
\midrule
CelebA & MSE & 0.640\statusfinal & 0.475\statusfinal & 0.499\statusfinal & 0.399\statusfinal \\
CelebA & SSIM & 0.828\statusfinal & 0.693\statusfinal & 0.844\statusfinal & 0.726\statusfinal \\
CelebA & ArcFace donor reference & 0.354\statusfinal & \textit{unavailable} & \textit{unavailable} & 0.983\statusfinal \\
IMDb-WIKI & MSE & 0.690\statusfinal & 0.529\statusfinal & 0.498\statusfinal & 0.397\statusfinal \\
IMDb-WIKI & SSIM & 0.856\statusfinal & 0.739\statusfinal & 0.859\statusfinal & 0.743\statusfinal \\
IMDb-WIKI & ArcFace donor reference & 0.396\statusfinal & \textit{unavailable} & \textit{unavailable} & 0.984\statusfinal \\
\bottomrule

\end{tabular}
\end{table*}

\subsection{Iris images (UTIRIS, UPOL)}
\label{sec:r-iris}
The current P2 strip panel contains 156 target eyes from 78 people for UTIRIS NIR.
The secondary UTIRIS VIS and UPOL panels contain 158 eyes from 79 people and 128 eyes from 64 people, respectively\statuspre.
The VIS repeat endpoint has one fewer eligible eye, which is retained as endpoint-specific coverage rather than silently changing every comparison.
Table~\ref{tab:iris} reports mean-relative squared error on the support defined for each condition.
The repeated eye is closer than the fellow eye in all three saved summaries, and a context-selected unrelated donor can fall between them.
These are pixel-space relationships; they are not an iris biometric verification result.

\begin{table*}[t]
\centering
\caption{\textbf{Provisional P2 iris donor ladder.} Squared-error ratio $Q$ against an independent dataset mean evaluated on the same support; lower is closer.
The mean anchor is one by construction; an independent equal-variance stranger has theoretical expectation near two.
Intervals use person bootstrap.
Target-eye/person counts are 156/78, 158/79, and 128/64; the VIS repeat endpoint uses 157 eyes.
Source: \src{results/iris_ladder_metrics.json}, results block.}
\label{tab:iris}
\small
\begin{tabular}{@{}lcccc@{}}
\toprule
Dataset & Same-eye repeat & Fellow eye & Context donor & Stranger \\
\midrule
UTIRIS NIR & \shortstack{0.407\statuspre\\{\scriptsize [0.345, 0.475]}} & \shortstack{1.256\statuspre\\{\scriptsize [1.123, 1.395]}} & \shortstack{0.599\statuspre\\{\scriptsize [0.540, 0.660]}} & \shortstack{2.001\statuspre\\{\scriptsize [1.816, 2.207]}} \\
UTIRIS VIS & \shortstack{0.569\statuspre\\{\scriptsize [0.503, 0.640]}} & \shortstack{1.293\statuspre\\{\scriptsize [1.211, 1.375]}} & \shortstack{0.861\statuspre\\{\scriptsize [0.826, 0.898]}} & \shortstack{2.018\statuspre\\{\scriptsize [1.884, 2.176]}} \\
UPOL & \shortstack{0.177\statuspre\\{\scriptsize [0.154, 0.200]}} & \shortstack{1.118\statuspre\\{\scriptsize [1.037, 1.198]}} & \shortstack{0.841\statuspre\\{\scriptsize [0.798, 0.882]}} & \shortstack{2.138\statuspre\\{\scriptsize [2.007, 2.281]}} \\
\bottomrule

\end{tabular}
\end{table*}

The source retains distinct $R_{\rm self}$ and $R_{\rm rep}$ values, window-level controls, and the directional tests.
It also records a lid-shadow mask review limitation.
The newer full P3 run is absent from this aggregate snapshot.
Consequently, its LPIPS, Betti-matching, FID, and FWD outcomes are not numerically reconstructed from P2 or from old iris-code results.

\subsection{Synthetic simulations}
\label{sec:r-sim}
The saved synthetic experiments isolate a complementary ambiguity.
On the known-target attenuation path, an image can keep a high unrelated-reference rank while its residual amplitude falls.
An amplitude-matched incorrect residual shows why distance from the mean alone is insufficient.
Synthetic pupil dilation holds latent eye parameters fixed while changing appearance, exposing the opposite failure direction.
These controls explain mechanisms, but they do not establish how often the same effect occurs in measured anatomy.
The synthetic label experiments additionally depend on the anatomical prior and segmentation construction.
Their conditional findings do not establish a universal Dice or topology result.
